\documentclass[12pt,a4paper]{article}

% ─── Paquetes ───────────────────────────────────────────────────────────────
\usepackage[utf8]{inputenc}
\usepackage[T1]{fontenc}
\usepackage[spanish]{babel}
\usepackage{geometry}
\usepackage{booktabs}
\usepackage{longtable}
\usepackage{array}
\usepackage{xcolor}
\usepackage{hyperref}
\usepackage{titlesec}
\usepackage{enumitem}
\usepackage{fancyhdr}
\usepackage{graphicx}
\usepackage{microtype}
\usepackage{parskip}
\usepackage{mdframed}
\usepackage{tocloft}

% ─── Geometría ───────────────────────────────────────────────────────────────
\geometry{
  top=2.5cm, bottom=2.5cm,
  left=3cm,  right=2.5cm
}

% ─── Colores ─────────────────────────────────────────────────────────────────
\definecolor{navyblue}{RGB}{0,51,102}
\definecolor{gpsblue}{RGB}{30,100,200}
\definecolor{alertorange}{RGB}{220,100,0}
\definecolor{notegreen}{RGB}{0,120,60}
\definecolor{warningred}{RGB}{180,0,0}
\definecolor{tablegray}{RGB}{245,245,245}
\definecolor{headerblue}{RGB}{0,70,130}

% ─── Hipervínculos ───────────────────────────────────────────────────────────
\hypersetup{
  colorlinks=true,
  linkcolor=navyblue,
  urlcolor=gpsblue,
  pdftitle={Manual de Usuario GEONAVAL},
  pdfauthor={GEONAVAL},
  pdfsubject={Sistema de Registro Fluvial y Monitoreo GPS}
}

% ─── Encabezados de sección ──────────────────────────────────────────────────
\titleformat{\section}[block]
  {\large\bfseries\color{navyblue}}
  {\thesection.}{0.6em}{}
  [\vspace{-4pt}\color{navyblue}\rule{\linewidth}{0.8pt}\vspace{4pt}]

\titleformat{\subsection}[hang]
  {\normalsize\bfseries\color{headerblue}}
  {\thesubsection}{0.6em}{}

% ─── Cabecera y pie de página ────────────────────────────────────────────────
\pagestyle{fancy}
\fancyhf{}
\fancyhead[L]{\small\textcolor{navyblue}{\textbf{GEONAVAL}} -- Manual de Usuario}
\fancyhead[R]{\small\textcolor{gray}{v1.0.0 -- Mayo 2026}}
\fancyfoot[C]{\small\thepage}
\renewcommand{\headrulewidth}{0.4pt}
\renewcommand{\footrulewidth}{0.2pt}

% ─── Entornos de nota / advertencia ─────────────────────────────────────────
\newmdenv[
  linecolor=notegreen,
  backgroundcolor=notegreen!8,
  leftmargin=0pt, rightmargin=0pt,
  innerleftmargin=8pt, innerrightmargin=8pt,
  innertopmargin=6pt, innerbottommargin=6pt,
  linewidth=3pt, topline=false, rightline=false, bottomline=false
]{notabox}

\newmdenv[
  linecolor=warningred,
  backgroundcolor=warningred!8,
  leftmargin=0pt, rightmargin=0pt,
  innerleftmargin=8pt, innerrightmargin=8pt,
  innertopmargin=6pt, innerbottommargin=6pt,
  linewidth=3pt, topline=false, rightline=false, bottomline=false
]{advertenciabox}

% ─── Separador de perfil ─────────────────────────────────────────────────────
\newcommand{\perfilbanner}[1]{%
  \vspace{1em}%
  \noindent\colorbox{navyblue}{\parbox{\linewidth}{\centering\large\bfseries\color{white}\strut #1\strut}}%
  \vspace{0.5em}%
}

% ─── Estilo de tablas ────────────────────────────────────────────────────────
\newcolumntype{L}{>{\raggedright\arraybackslash}p{4cm}}
\newcolumntype{R}{>{\raggedright\arraybackslash}p{10cm}}

\setlength{\arrayrulewidth}{0.4pt}
\arrayrulecolor{gray!60}

% ════════════════════════════════════════════════════════════════════════════
\begin{document}

% ─── PORTADA ─────────────────────────────────────────────────────────────────
\begin{titlepage}
  \centering
  \vspace*{2cm}

  {\Huge\bfseries\color{navyblue} GEONAVAL\par}
  \vspace{0.4cm}
  {\Large\color{gpsblue} Control Fluvial\par}
  \vspace{1.2cm}
  {\huge\bfseries Manual de Usuario\par}
  \vspace{0.6cm}
  {\large\itshape Sistema de Registro Fluvial y Monitoreo GPS\par}
  \vspace{0.4cm}
  {\normalsize Control y Seguridad del Transporte Fluvial\par}
  \vspace{0.4cm}
  {\small Incluye perfil Operador · Perfil Administrador · Centro de Auditoría\par}

  \vspace{2cm}
  \begin{tabular}{ll}
    \toprule
    \textbf{Versión}      & 1.0.0 -- Mayo 2026 \\
    \textbf{Materia}      & Ingeniería de Software \\
    \textbf{Ubicación}    & Quibdó, Chocó -- Colombia \\
    \textbf{Fecha}        & Mayo 2026 \\
    \textbf{URL del sistema} & \url{https://geonaval.onrender.com} \\
    \bottomrule
  \end{tabular}

  \vfill
  {\small Quibdó, Chocó -- Colombia \quad|\quad Versión 1.0.0 -- Mayo 2026\par}
\end{titlepage}

% ─── TABLA DE CONTENIDO ──────────────────────────────────────────────────────
\tableofcontents
\newpage

% ════════════════════════════════════════════════════════════════════════════
\section{Introducción}

\subsection{Descripción del Sistema}

GEONAVAL es un sistema web de Registro Fluvial y Monitoreo GPS diseñado para el control y la seguridad del transporte fluvial en la región de Quibdó, Chocó, Colombia. El sistema permite gestionar de forma integral todos los actores y procesos involucrados en la operación fluvial, incluyendo embarcaciones, propietarios, tripulación, pasajeros y viajes.

El sistema cuenta con monitoreo GPS en tiempo real, permitiendo a las autoridades y administradores hacer seguimiento a las embarcaciones activas sobre las rutas fluviales del río Atrato y sus afluentes.

\subsection{Propósito del Manual}

Este manual tiene como objetivo guiar a todos los perfiles de usuario del sistema GEONAVAL --- Operador, Administrador y Autoridad --- en el uso correcto de todas las funcionalidades disponibles. Incluye capturas de pantalla reales del sistema con descripciones detalladas de cada elemento.

\subsection{Alcance}

El presente manual cubre los siguientes módulos del sistema:

\begin{itemize}[leftmargin=1.5em]
  \item Inicio de sesión y autenticación
  \item Perfil Operador: Mis Viajes, Mis Pasajeros, Mi Ruta GPS, Reportar Incidente, Notificaciones, Configuración
  \item Perfil Administrador: Dashboard, Embarcaciones, Propietarios, Tripulación, Compras y Tickets
  \item Monitoreo GPS en Tiempo Real
  \item Reportes y Estadísticas
  \item Centro de Notificaciones y Auditoría del Sistema
  \item Configuración del Sistema
\end{itemize}

\subsection{Convenciones del Manual}

\begin{longtable}{LR}
  \toprule
  \textbf{Elemento} & \textbf{Descripción} \\
  \midrule
  \endfirsthead
  \toprule
  \textbf{Elemento} & \textbf{Descripción} \\
  \midrule
  \endhead
  \bottomrule
  \endfoot
  Texto en \textbf{negrita}   & Indica nombres de botones, campos o elementos de interfaz. \\[4pt]
  {[}Texto entre corchetes{]} & Indica una acción que debe realizar el usuario. \\[4pt]
  \textbf{Nota:}              & Información importante a tener en cuenta. \\[4pt]
  \textbf{Advertencia:}       & Acción que puede tener consecuencias irreversibles. \\[4pt]
  Flechas rojas ($\rightarrow$) & Señalan los elementos clave en cada captura de pantalla. \\[4pt]
  ①②③\ldots               & Numerales en círculo que señalan elementos en las capturas de pantalla. \\
\end{longtable}

% ════════════════════════════════════════════════════════════════════════════
\section{Requisitos del Sistema}

\subsection{Requisitos Técnicos}

\begin{longtable}{LR}
  \toprule
  \textbf{Campo} & \textbf{Descripción} \\
  \midrule
  \endfirsthead
  \toprule
  \textbf{Campo} & \textbf{Descripción} \\
  \midrule
  \endhead
  \bottomrule
  \endfoot
  Navegador web        & Google Chrome 90+, Mozilla Firefox 88+, Microsoft Edge 90+, Safari 14+ \\[4pt]
  Conexión a Internet  & Mínimo 2 Mbps (recomendado 5 Mbps o superior para GPS en tiempo real) \\[4pt]
  Resolución de pantalla & Mínimo 1280$\times$720 px (recomendado 1920$\times$1080) \\[4pt]
  Sistema operativo    & Windows 10+, macOS 10.14+, Ubuntu 18.04+ \\[4pt]
  JavaScript           & Habilitado en el navegador \\
\end{longtable}

\subsection{Acceso al Sistema}

GEONAVAL es una aplicación web; no requiere instalación. El acceso se realiza directamente desde el navegador ingresando la URL del sistema:

\begin{center}
  \large\url{https://geonaval.onrender.com}
\end{center}

% ════════════════════════════════════════════════════════════════════════════
\section{Inicio de Sesión}

\subsection{Acceso al Sistema}

Para ingresar a GEONAVAL, el usuario debe completar el formulario de autenticación que aparece al abrir el sistema en el navegador.

\textit{Figura 1 -- Pantalla de Inicio de Sesión de GEONAVAL}

\subsection{Descripción de Elementos (ver Figura 1)}

\begin{longtable}{LR}
  \toprule
  \textbf{Campo} & \textbf{Descripción} \\
  \midrule
  \endfirsthead
  \toprule
  \textbf{Campo} & \textbf{Descripción} \\
  \midrule
  \endhead
  \bottomrule
  \endfoot
  ① Correo Electrónico & Campo donde se ingresa el correo registrado del usuario (ejemplo: admin@geonaval.com). \\[4pt]
  ② Contraseña         & Campo de contraseña. El ícono de ojo permite mostrar u ocultar el texto. \\[4pt]
  ③ Tipo de Usuario    & Menú desplegable para seleccionar el perfil: Administrador, Operador o Autoridad. \\[4pt]
  ④ Iniciar Sesión     & Botón principal para autenticarse en el sistema. \\[4pt]
  ⑤ Recuperar contraseña & Enlace para restablecer la contraseña olvidada vía correo electrónico. \\
\end{longtable}

\subsection{Pasos para Iniciar Sesión}

\begin{enumerate}[leftmargin=1.5em]
  \item Abra su navegador web preferido.
  \item Ingrese la URL del sistema GEONAVAL en la barra de direcciones.
  \item Complete el campo \textbf{Correo Electrónico} (①).
  \item Complete el campo \textbf{Contraseña} (②).
  \item Seleccione su \textbf{Tipo de Usuario} en el desplegable (③).
  \item Haga clic en \textbf{Iniciar Sesión} (④).
\end{enumerate}

\subsection{Tipos de Usuario}

\begin{longtable}{LR}
  \toprule
  \textbf{Perfil} & \textbf{Descripción} \\
  \midrule
  \endfirsthead
  \toprule
  \textbf{Perfil} & \textbf{Descripción} \\
  \midrule
  \endhead
  \bottomrule
  \endfoot
  Administrador & Acceso completo a todos los módulos del sistema. \\[4pt]
  Operador      & Acceso a módulos operativos: viajes, pasajeros, GPS, incidentes y notificaciones. \\[4pt]
  Autoridad     & Acceso de solo lectura a monitoreo, reportes y alertas. \\
\end{longtable}

% ════════════════════════════════════════════════════════════════════════════
\perfilbanner{PERFIL: OPERADOR}

\section{Perfil Operador -- Mis Viajes}

\subsection{Descripción General}

La pantalla \textbf{Mis Viajes} es el panel de control principal del Operador. Muestra un resumen del estado operativo, la lista de pasajeros registrados y los próximos viajes programados.

\textit{Figura 1 -- Pantalla Mis Viajes (Perfil Operador)}

\subsection{Indicadores KPI}

\begin{longtable}{LR}
  \toprule
  \textbf{KPI} & \textbf{Descripción} \\
  \midrule
  \endfirsthead
  \toprule
  \textbf{KPI} & \textbf{Descripción} \\
  \midrule
  \endhead
  \bottomrule
  \endfoot
  Viajes Asignados & Número de viajes asignados al operador activo. En estado inicial muestra 0. \\[4pt]
  Viaje en Curso   & Número de viajes actualmente en ejecución. El administrador debe iniciarlos. \\[4pt]
  Pasajeros        & Total de pasajeros registrados por el operador en el sistema. \\[4pt]
  Llegada Est.     & Hora estimada de llegada al destino del viaje en curso. \\
\end{longtable}

\subsection{Descripción de Elementos (Figura 1 -- Mis Viajes)}

\begin{longtable}{LR}
  \toprule
  \textbf{Campo} & \textbf{Descripción} \\
  \midrule
  \endfirsthead
  \toprule
  \textbf{Campo} & \textbf{Descripción} \\
  \midrule
  \endhead
  \bottomrule
  \endfoot
  ① Bienvenido, Operador          & Encabezado de bienvenida que muestra el nombre del perfil activo. \\[4pt]
  ② Sin embarcación / 0 viaje(s)  & Estado actual de la embarcación asignada. \\[4pt]
  ③ KPIs superiores               & Cuatro indicadores: Viajes Asignados, Viaje en Curso, Pasajeros y Llegada Est. \\[4pt]
  ④ Mensaje informativo           & Indica si no hay viaje en curso; el administrador debe iniciar uno. \\[4pt]
  ⑤ Tabla Pasajeros registrados   & Lista todos los pasajeros con: Nombre, Documento, Destino y Estado. \\[4pt]
  ⑥ Próximos Viajes Programados   & Panel con los viajes programados pendientes de inicio. \\[4pt]
  ⑦ Reportar Incidente            & Acceso rápido al formulario de reporte de incidentes desde el dashboard. \\
\end{longtable}

\subsection{Tabla de Pasajeros Registrados}

\begin{longtable}{LLLL}
  \toprule
  \textbf{Nombre} & \textbf{Documento} & \textbf{Destino/Ruta} & \textbf{Estado} \\
  \midrule
  \endfirsthead
  \toprule
  \textbf{Nombre} & \textbf{Documento} & \textbf{Destino/Ruta} & \textbf{Estado} \\
  \midrule
  \endhead
  \bottomrule
  \endfoot
  Nombre completo & N.° de identificación & Ciudad o punto fluvial & Confirmado / Pendiente \\
\end{longtable}

\subsection{Estado del Viaje}

\begin{notabox}
\textbf{Nota:} Para que el operador pueda iniciar el tracking GPS y ver viajes en curso, el Administrador debe cambiar el estado del viaje a \textbf{`En Curso'} desde el panel de administración.
\end{notabox}

% ════════════════════════════════════════════════════════════════════════════
\section{Perfil Operador -- Mis Pasajeros}

\subsection{Descripción}

El módulo \textbf{Mis Pasajeros} permite al operador consultar el listado completo de pasajeros registrados para sus viajes asignados.

\textit{Figura 2 -- Módulo Mis Pasajeros (Perfil Operador)}

\subsection{Descripción de Elementos (Figura 2 -- Mis Pasajeros)}

\begin{longtable}{LR}
  \toprule
  \textbf{Campo} & \textbf{Descripción} \\
  \midrule
  \endfirsthead
  \toprule
  \textbf{Campo} & \textbf{Descripción} \\
  \midrule
  \endhead
  \bottomrule
  \endfoot
  ① Encabezado del módulo & Muestra el título `Mis Pasajeros' y el conteo total de pasajeros. \\[4pt]
  ② Columna NOMBRE        & Nombre completo del pasajero. \\[4pt]
  ③ Columna DOCUMENTO     & Número de identificación del pasajero. \\[4pt]
  ④ Columna DESTINO/RUTA  & Nombre del destino fluvial asignado al pasajero. \\[4pt]
  ⑤ Columna ESTADO        & Estado del tiquete: Confirmado (azul) / Pendiente (amarillo). \\[4pt]
  ⑥ Barra de desplazamiento & Permite navegar verticalmente si hay más pasajeros. \\
\end{longtable}

\subsection{Listado de Pasajeros -- Ejemplo}

\begin{longtable}{llll}
  \toprule
  \textbf{Nombre} & \textbf{Documento} & \textbf{Destino/Ruta} & \textbf{Estado} \\
  \midrule
  \endfirsthead
  \toprule
  \textbf{Nombre} & \textbf{Documento} & \textbf{Destino/Ruta} & \textbf{Estado} \\
  \midrule
  \endhead
  \bottomrule
  \endfoot
  ddddddd         & 11111111   & Tadó    & Confirmado \\
  qweqweqeeee     & 2333244    & Tadó    & Confirmado \\
  erwrwerwer      & 341441421  & Tadó    & Confirmado \\
  popopopopop     & 1010101010 & Tadó    & Confirmado \\
  sdadasd         & 12313213   & Istmina & Confirmado \\
\end{longtable}

% ════════════════════════════════════════════════════════════════════════════
\section{Perfil Operador -- Mi Ruta GPS}

\subsection{Descripción}

El módulo \textbf{Mi Ruta GPS} permite al operador compartir su ubicación en tiempo real durante un viaje activo. El mapa muestra las embarcaciones en tránsito sobre el Río Atrato, Quibdó.

\textit{Figura 3 -- Módulo Mi Ruta GPS (Perfil Operador)}

\subsection{Descripción de Elementos (Figura 3 -- Mi Ruta GPS)}

\begin{longtable}{LR}
  \toprule
  \textbf{Campo} & \textbf{Descripción} \\
  \midrule
  \endfirsthead
  \toprule
  \textbf{Campo} & \textbf{Descripción} \\
  \midrule
  \endhead
  \bottomrule
  \endfoot
  ① Selector de viaje en curso & Desplegable para seleccionar el viaje activo antes de iniciar el tracking. \\[4pt]
  ② Botón Iniciar Tracking     & Activa el envío de ubicación GPS en tiempo real. Solo disponible si hay viaje en curso. \\[4pt]
  ③ Mensaje de advertencia     & Indica que no hay viajes en curso y que el administrador debe cambiar el estado. \\[4pt]
  ④ Mapa -- Río Atrato         & Mapa interactivo (OpenStreetMap + Leaflet) centrado en Quibdó. Actualización cada 5 s. \\[4pt]
  ⑤ Indicador EN VIVO          & Punto rojo que confirma actualizaciones en tiempo real. \\[4pt]
  ⑥ Leyenda del mapa           & Guía de colores: azul = embarcación en curso, verde = operador, línea = recorrido. \\[4pt]
  ⑦ Embarcaciones en Tránsito  & Panel con la lista de embarcaciones activas en el momento. \\[4pt]
  ⑧ Estado del Sistema GPS     & Tabla con estado del servidor GPS, mapa, viajes monitoreados y frecuencia. \\
\end{longtable}

\subsection{Estado del Sistema GPS}

\begin{longtable}{LR}
  \toprule
  \textbf{Campo} & \textbf{Descripción} \\
  \midrule
  \endfirsthead
  \toprule
  \textbf{Campo} & \textbf{Descripción} \\
  \midrule
  \endhead
  \bottomrule
  \endfoot
  Servidor GPS       & Activo -- el servidor de posicionamiento está en línea. \\[4pt]
  Mapa OpenStreetMap & Conectado -- el mapa de fondo está cargado correctamente. \\[4pt]
  Viajes monitoreados & Número de viajes activos rastreados en tiempo real. \\[4pt]
  Actualización      & Cada 5 segundos -- frecuencia de refresco de posición. \\
\end{longtable}

\begin{notabox}
\textbf{Nota:} Para que aparezca la posición en el mapa, el operador debe seleccionar un viaje en curso en el desplegable y hacer clic en \textbf{Iniciar Tracking}. La ubicación se obtendrá del dispositivo GPS del operador o del navegador del celular.
\end{notabox}

% ════════════════════════════════════════════════════════════════════════════
\section{Perfil Operador -- Reportar Incidente}

\subsection{Descripción}

El módulo \textbf{Reportar Incidente} permite al operador registrar formalmente cualquier situación anormal ocurrida durante la operación fluvial: retrasos, averías, emergencias u otras situaciones.

\textit{Figura 4 -- Módulo Reportar Incidente (Perfil Operador)}

\subsection{Descripción de Elementos (Figura 4 -- Reportar Incidente)}

\begin{longtable}{LR}
  \toprule
  \textbf{Campo} & \textbf{Descripción} \\
  \midrule
  \endfirsthead
  \toprule
  \textbf{Campo} & \textbf{Descripción} \\
  \midrule
  \endhead
  \bottomrule
  \endfoot
  ① Viaje relacionado (opcional) & Desplegable para asociar el incidente a un viaje. Si no aplica, seleccionar `Sin viaje específico'. \\[4pt]
  ② Tipo de incidente            & Categoría: Retraso, Avería, Emergencia u Otro. \\[4pt]
  ③ Severidad                   & Nivel de gravedad: Baja, Media, Alta o Crítica. \\[4pt]
  ④ Descripción                 & Campo de texto libre para describir el incidente. \\[4pt]
  ⑤ Enviar reporte              & Botón naranja que envía el reporte y notifica al Administrador. \\
\end{longtable}

\subsection{Pasos para Reportar un Incidente}

\begin{enumerate}[leftmargin=1.5em]
  \item En el menú lateral, haga clic en \textbf{Reportar Incidente}.
  \item Seleccione el viaje relacionado (opcional).
  \item Elija el \textbf{Tipo de incidente} en el desplegable (②).
  \item Seleccione el nivel de \textbf{Severidad} (③).
  \item Escriba una descripción detallada en el campo de texto (④).
  \item Haga clic en \textbf{Enviar reporte} (⑤) para registrar el incidente.
\end{enumerate}

\begin{advertenciabox}
\textbf{⚠ Advertencia:} Una vez enviado, el reporte queda registrado en el sistema y es visible para el Administrador. No es posible eliminar un reporte enviado.
\end{advertenciabox}

% ════════════════════════════════════════════════════════════════════════════
\section{Perfil Operador -- Notificaciones y Auditoría del Sistema}

\subsection{Centro de Notificaciones}

El módulo de \textbf{Notificaciones} muestra el historial de todos los eventos y cambios registrados en el sistema. Funciona también como Centro de Auditoría para el Administrador.

\textit{Figura 5 -- Centro de Notificaciones y Auditoría del Sistema}

\subsection{Estadísticas del Centro (Figura 5)}

\begin{longtable}{LR}
  \toprule
  \textbf{Campo} & \textbf{Descripción} \\
  \midrule
  \endfirsthead
  \toprule
  \textbf{Campo} & \textbf{Descripción} \\
  \midrule
  \endhead
  \bottomrule
  \endfoot
  Total     & Número total de notificaciones/eventos registrados. \\[4pt]
  Leídas    & Notificaciones que ya han sido revisadas. \\[4pt]
  Incidentes & Cantidad de incidentes reportados. \\[4pt]
  Recientes & Notificaciones generadas en las últimas horas. \\
\end{longtable}

\subsection{Tipos de Eventos Registrados}

\begin{longtable}{LR}
  \toprule
  \textbf{Tipo} & \textbf{Descripción} \\
  \midrule
  \endfirsthead
  \toprule
  \textbf{Tipo} & \textbf{Descripción} \\
  \midrule
  \endhead
  \bottomrule
  \endfoot
  VIAJE (azul)          & Cambios de estado de viajes: Cancelado, Iniciado, Finalizado, Programado. \\[4pt]
  PASAJERO (verde)      & Registro, modificación o inscripción de pasajeros en viajes. \\[4pt]
  EMBARCACIÓN (naranja) & Registro de nuevas embarcaciones en el sistema. \\[4pt]
  TRIPULACIÓN (morado)  & Registro de nuevos miembros de tripulación. \\[4pt]
  USUARIO (gris)        & Creación, modificación o eliminación de cuentas de usuario. \\[4pt]
  COMPRA (amarillo)     & Compra de pasajes por parte de pasajeros. \\
\end{longtable}

\subsection{Filtros Disponibles}

La barra superior permite filtrar notificaciones por categoría: Todos, Viaje, Embarcaciones, Tripulación, Pasajeros, Propietarios, Usuarios, Incidentes.

\subsection{Botones de Acción}

\begin{longtable}{LR}
  \toprule
  \textbf{Botón} & \textbf{Descripción} \\
  \midrule
  \endfirsthead
  \toprule
  \textbf{Botón} & \textbf{Descripción} \\
  \midrule
  \endhead
  \bottomrule
  \endfoot
  Actualizar               & Recarga la lista con los eventos más recientes del servidor. \\[4pt]
  Marcar todos como leídos & Marca todas las notificaciones como revisadas de forma masiva. \\
\end{longtable}

% ════════════════════════════════════════════════════════════════════════════
\section{Perfil Operador -- Configuración}

\subsection{Descripción}

El módulo de \textbf{Configuración} permite al operador ver y gestionar los datos de su perfil, opciones de seguridad y la información de la versión del sistema.

\textit{Figura 6 -- Módulo Configuración (Perfil Operador)}

\subsection{Descripción de Elementos (Figura 6 -- Configuración)}

\begin{longtable}{LR}
  \toprule
  \textbf{Campo} & \textbf{Descripción} \\
  \midrule
  \endfirsthead
  \toprule
  \textbf{Campo} & \textbf{Descripción} \\
  \midrule
  \endhead
  \bottomrule
  \endfoot
  ① Editar Perfil             & Botón para modificar los datos del usuario activo. \\[4pt]
  ② Datos del usuario activo  & Nombre completo, correo, rol, último acceso y fecha de registro. \\[4pt]
  ③ Sección Seguridad         & Incluye: cambiar contraseña e historial de sesiones. \\[4pt]
  Cambiar contraseña          & Permite actualizar la contraseña de acceso. \\[4pt]
  Ver historial de sesiones   & Últimas sesiones iniciadas con fecha, hora y dispositivo. \\[4pt]
  ④ Información del sistema   & Versión del sistema (1.0.0) y fecha de última actualización. \\
\end{longtable}

\subsection{Pasos para Cambiar la Contraseña}

\begin{enumerate}[leftmargin=1.5em]
  \item Haga clic en \textbf{Configuración} en el menú lateral.
  \item En la sección \textbf{Seguridad}, haga clic en \textbf{Cambiar contraseña}.
  \item Ingrese la contraseña actual.
  \item Ingrese la nueva contraseña y confírmela.
  \item Haga clic en \textbf{Guardar cambios}.
\end{enumerate}

% ════════════════════════════════════════════════════════════════════════════
\perfilbanner{PERFIL: ADMINISTRADOR}

\section{Perfil Administrador -- Dashboard Principal}

\subsection{Descripción General}

Una vez autenticado como Administrador, el sistema redirige al \textbf{Dashboard Principal}. Este panel muestra los indicadores clave de la operación fluvial y acceso directo al mapa GPS en tiempo real.

\textit{Figura 7 -- Dashboard Principal de GEONAVAL}

\subsection{Descripción de Elementos (ver Figura 7)}

\begin{longtable}{LR}
  \toprule
  \textbf{Campo} & \textbf{Descripción} \\
  \midrule
  \endfirsthead
  \toprule
  \textbf{Campo} & \textbf{Descripción} \\
  \midrule
  \endhead
  \bottomrule
  \endfoot
  ① Menú de navegación lateral & Panel izquierdo con acceso a todos los módulos del sistema. \\[4pt]
  ② Embarcaciones Activas      & KPI con el número de embarcaciones operativas y variación mensual. \\[4pt]
  ③ Viajes en Curso            & KPI con el número de viajes activos y variación mensual. \\[4pt]
  ④ Perfil de usuario          & Área superior derecha con nombre, rol y opciones de sesión. \\[4pt]
  ⑤ Mapa GPS en tiempo real   & Vista del mapa fluvial con posición de embarcaciones activas. \\[4pt]
  ⑥ Próxima Llegada           & Panel con hora estimada y datos de la siguiente embarcación. \\
\end{longtable}

\subsection{Indicadores Clave (KPIs)}

Los cuatro KPIs de la parte superior muestran: \textbf{Embarcaciones Activas}, \textbf{Viajes en Curso}, \textbf{Pasajeros Transportados} y \textbf{Alertas Activas}, todos con comparativo frente al mes anterior.

% ════════════════════════════════════════════════════════════════════════════
\section{Perfil Administrador -- Gestión de Embarcaciones}

El módulo de \textbf{Gestión de Embarcaciones} permite registrar, consultar y administrar todas las embarcaciones fluviales.

\textit{Figura 8 -- Módulo de Gestión de Embarcaciones}

\subsection{Descripción de Elementos (ver Figura 8)}

\begin{longtable}{LR}
  \toprule
  \textbf{Campo} & \textbf{Descripción} \\
  \midrule
  \endfirsthead
  \toprule
  \textbf{Campo} & \textbf{Descripción} \\
  \midrule
  \endhead
  \bottomrule
  \endfoot
  ① Nueva Embarcación      & Botón para abrir el formulario de registro de una nueva embarcación. \\[4pt]
  ② KPIs del módulo        & Resumen de totales: registradas, operativas, en mantenimiento y viajes. \\[4pt]
  ③ Estado: Operativa      & Etiqueta verde --- embarcación disponible para viajes. \\[4pt]
  ④ Estado: En Mantenimiento & Etiqueta naranja --- embarcación temporalmente fuera de servicio. \\[4pt]
  ⑤ Ver detalles           & Enlace para consultar toda la información de una embarcación. \\
\end{longtable}

\subsection{Registrar Nueva Embarcación}

\begin{enumerate}[leftmargin=1.5em]
  \item Haga clic en \textbf{+ Nueva Embarcación} (①).
  \item Complete el formulario: nombre, tipo, propietario y estado inicial.
  \item Haga clic en \textbf{Guardar} para registrar.
\end{enumerate}

% ════════════════════════════════════════════════════════════════════════════
\section{Perfil Administrador -- Gestión de Propietarios}

Este módulo permite administrar el registro de propietarios de embarcaciones, tanto personas naturales como empresas.

\textit{Figura 9 -- Módulo de Gestión de Propietarios}

\subsection{Descripción de Elementos (ver Figura 9)}

\begin{longtable}{LR}
  \toprule
  \textbf{Campo} & \textbf{Descripción} \\
  \midrule
  \endfirsthead
  \toprule
  \textbf{Campo} & \textbf{Descripción} \\
  \midrule
  \endhead
  \bottomrule
  \endfoot
  ① Nuevo Propietario         & Botón para registrar un nuevo propietario en el sistema. \\[4pt]
  ② Tipo: Natural / Empresa   & Clasificación del propietario como persona física o jurídica. \\[4pt]
  ③ Ver / Editar / Eliminar   & Íconos de acción disponibles por cada registro. \\[4pt]
  ④ Embarcaciones asociadas   & Contador que muestra cuántas embarcaciones tiene cada propietario. \\
\end{longtable}

\begin{advertenciabox}
\textbf{⚠ Advertencia:} Al eliminar un propietario esta acción es irreversible. Verifique que no tenga embarcaciones activas antes de proceder.
\end{advertenciabox}

% ════════════════════════════════════════════════════════════════════════════
\section{Perfil Administrador -- Gestión de Tripulación}

El módulo de \textbf{Tripulación} permite administrar el personal operativo de las embarcaciones fluviales.

\textit{Figura 10 -- Módulo de Gestión de Tripulación}

\subsection{Descripción de Elementos (ver Figura 10)}

\begin{longtable}{LR}
  \toprule
  \textbf{Campo} & \textbf{Descripción} \\
  \midrule
  \endfirsthead
  \toprule
  \textbf{Campo} & \textbf{Descripción} \\
  \midrule
  \endhead
  \bottomrule
  \endfoot
  ① Nuevo Tripulante   & Botón para registrar un nuevo miembro de la tripulación. \\[4pt]
  ② KPIs del módulo    & Total de tripulantes, operadores fluviales, personal auxiliar y activos. \\[4pt]
  ③ Rol del tripulante & Cargo asignado: Operador Fluvial, Auxiliar de Pasajeros, Motorista, etc. \\[4pt]
  ④ Estado: Activo     & Etiqueta verde que indica disponibilidad del tripulante. \\[4pt]
  ⑤ Acciones          & Íconos para Ver, Editar y Eliminar cada registro. \\
\end{longtable}

\subsection{Registrar Nuevo Tripulante}

\begin{enumerate}[leftmargin=1.5em]
  \item Haga clic en \textbf{+ Nuevo Tripulante} (①).
  \item Complete nombre completo, rol, embarcación asignada y horario.
  \item Defina el estado inicial (Activo/Inactivo) y haga clic en \textbf{Guardar}.
\end{enumerate}

% ════════════════════════════════════════════════════════════════════════════
\section{Perfil Administrador -- Compras y Tickets}

Este módulo centraliza la venta de pasajes y la gestión de tickets de transporte fluvial.

\textit{Figura 11 -- Módulo de Compras y Tickets}

\subsection{Descripción de Elementos (ver Figura 11)}

\begin{longtable}{LR}
  \toprule
  \textbf{Campo} & \textbf{Descripción} \\
  \midrule
  \endfirsthead
  \toprule
  \textbf{Campo} & \textbf{Descripción} \\
  \midrule
  \endhead
  \bottomrule
  \endfoot
  ① Nueva Compra         & Botón para registrar una nueva venta de ticket. \\[4pt]
  ② KPIs del módulo      & Ventas del día, total recaudado, tickets confirmados y pendientes. \\[4pt]
  ③ Barra de búsqueda    & Permite buscar tickets por número, pasajero o documento. \\[4pt]
  ④ Filtrar por Fecha    & Abre un selector de rango de fechas para filtrar los registros. \\[4pt]
  ⑤ Exportar            & Descarga el listado de tickets en formato Excel. \\[4pt]
  ⑥ Estado Confirmado   & Etiqueta azul que indica que el pago ha sido verificado. \\[4pt]
  ⑦ Estado Pendiente    & Etiqueta amarilla que indica que el ticket aún requiere confirmación. \\
\end{longtable}

\subsection{Crear Nueva Compra}

\begin{enumerate}[leftmargin=1.5em]
  \item Haga clic en \textbf{+ Nueva Compra} (①).
  \item Seleccione el pasajero o regístrelo si es nuevo.
  \item Elija la ruta, fecha, asiento y método de pago.
  \item Confirme la transacción para generar el ticket.
\end{enumerate}

% ════════════════════════════════════════════════════════════════════════════
\section{Perfil Administrador -- Monitoreo GPS en Tiempo Real}

El módulo de \textbf{Monitoreo GPS} permite hacer seguimiento en tiempo real de las embarcaciones activas sobre el Río Atrato y sus afluentes.

\textit{Figura 12 -- Monitoreo GPS en Tiempo Real}

\subsection{Descripción de Elementos (ver Figura 12)}

\begin{longtable}{LR}
  \toprule
  \textbf{Campo} & \textbf{Descripción} \\
  \midrule
  \endfirsthead
  \toprule
  \textbf{Campo} & \textbf{Descripción} \\
  \midrule
  \endhead
  \bottomrule
  \endfoot
  ① Indicador EN VIVO           & Punto rojo parpadeante que confirma datos en tiempo real. \\[4pt]
  ② Embarcación en curso (verde) & Punto verde que representa una embarcación navegando actualmente. \\[4pt]
  ③ Embarcación programada (azul) & Punto azul que indica un viaje confirmado pendiente de inicio. \\[4pt]
  ④ Leyenda del mapa            & Guía visual: verde = en curso, azul = programado, naranja = alerta. \\
\end{longtable}

\begin{notabox}
\textbf{Nota:} Si una embarcación no aparece en el mapa, es posible que su dispositivo GPS esté apagado o sin señal. Verifique el dispositivo a bordo.
\end{notabox}

% ════════════════════════════════════════════════════════════════════════════
\section{Perfil Administrador -- Reportes y Estadísticas}

El módulo de \textbf{Reportes} ofrece herramientas de análisis estadístico para evaluar la operación fluvial por períodos.

\subsection{Filtros y Gráficas}

\textit{Figura 13 -- Reportes: Filtros de Búsqueda y Gráficas Estadísticas}

\begin{longtable}{LR}
  \toprule
  \textbf{Campo} & \textbf{Descripción} \\
  \midrule
  \endfirsthead
  \toprule
  \textbf{Campo} & \textbf{Descripción} \\
  \midrule
  \endhead
  \bottomrule
  \endfoot
  ① Filtros de búsqueda  & Fecha Inicio, Fecha Fin, Embarcación, Ruta y Estado para acotar el reporte. \\[4pt]
  ② Aplicar / Limpiar    & Botones para aplicar los filtros o restablecer valores por defecto. \\[4pt]
  ③ Viajes/Pasajeros     & Gráfica de líneas con evolución mensual de viajes y pasajeros. \\[4pt]
  ④ Uso de Embarcaciones & Gráfica de barras con viajes y horas de operación por embarcación. \\[4pt]
  ⑤ Tooltip interactivo  & Al pasar el cursor sobre la gráfica se muestran los valores exactos del mes. \\
\end{longtable}

\subsection{Descarga de Reportes}

\textit{Figura 14 -- Reportes: Tipos de Reportes y Descarga}

\begin{longtable}{LR}
  \toprule
  \textbf{Campo} & \textbf{Descripción} \\
  \midrule
  \endfirsthead
  \toprule
  \textbf{Campo} & \textbf{Descripción} \\
  \midrule
  \endhead
  \bottomrule
  \endfoot
  ① Tipo de reporte      & Nombre y descripción del reporte (Viajes del Mes, Pasajeros, Rutas, etc.). \\[4pt]
  ② Descargar PDF        & Descarga el reporte en PDF (solo lectura, ideal para compartir). \\[4pt]
  ③ Descargar Excel      & Descarga el reporte en Excel (editable, ideal para análisis). \\[4pt]
  ④ Resumen estadístico  & Total viajes, pasajeros, horas y tasa de finalización del período. \\[4pt]
  ⑤ Pasajeros totales    & Cifra total de pasajeros transportados en el período seleccionado. \\
\end{longtable}

% ════════════════════════════════════════════════════════════════════════════
\section{Perfil Administrador -- Configuración del Sistema}

El módulo de \textbf{Configuración} permite gestionar el perfil, notificaciones y seguridad del sistema desde la vista de Administrador.

\textit{Figura 15 -- Módulo de Configuración del Sistema}

\subsection{Descripción de Elementos (ver Figura 15)}

\begin{longtable}{LR}
  \toprule
  \textbf{Campo} & \textbf{Descripción} \\
  \midrule
  \endfirsthead
  \toprule
  \textbf{Campo} & \textbf{Descripción} \\
  \midrule
  \endhead
  \bottomrule
  \endfoot
  ① Editar Perfil              & Botón para modificar los datos del usuario activo. \\[4pt]
  ② Datos del usuario activo   & Nombre, correo, rol, último acceso y fecha de registro. \\[4pt]
  ③ Notificaciones activas (✓) & Casillas marcadas para alertas de viajes, GPS y embarcaciones. \\[4pt]
  ④ Notificación desactivada   & Casilla desmarcada para reportes automáticos (desactivados por defecto). \\[4pt]
  ⑤ Opciones de Seguridad      & Accesos a: cambiar contraseña, 2FA, dispositivos e historial de sesiones. \\
\end{longtable}

% ════════════════════════════════════════════════════════════════════════════
\section{Glosario de Términos}

\begin{longtable}{LR}
  \toprule
  \textbf{Término} & \textbf{Definición} \\
  \midrule
  \endfirsthead
  \toprule
  \textbf{Término} & \textbf{Definición} \\
  \midrule
  \endhead
  \bottomrule
  \endfoot
  Dashboard           & Pantalla principal con indicadores y resúmenes visuales de la operación. \\[4pt]
  Embarcación         & Cualquier nave fluvial registrada: ferry, lancha, bote, etc. \\[4pt]
  GPS                 & Sistema de Posicionamiento Global. Permite rastrear la ubicación de las embarcaciones. \\[4pt]
  KPI                 & Indicador Clave de Rendimiento. Métrica para evaluar el desempeño del sistema. \\[4pt]
  Monitoreo en tiempo real & Seguimiento continuo y actualizado de posición y estado de embarcaciones. \\[4pt]
  Propietario         & Persona natural o jurídica dueña de una o más embarcaciones registradas. \\[4pt]
  Ticket              & Documento digital que representa el derecho de un pasajero a viajar. \\[4pt]
  Tripulación         & Personal operativo de una embarcación: operadores, auxiliares, motoristas. \\[4pt]
  Zarpe               & Autorización formal para que una embarcación parta hacia su destino. \\[4pt]
  Ruta fluvial        & Trayecto definido por el río que une dos puntos (origen y destino). \\[4pt]
  NIT                 & Número de Identificación Tributaria. Código fiscal de empresas en Colombia. \\[4pt]
  Alerta              & Notificación ante una situación anormal: desviación de ruta, fallo GPS, etc. \\[4pt]
  Tracking            & Proceso de envío continuo de coordenadas GPS desde el dispositivo del operador. \\[4pt]
  Incidente           & Evento registrado por el operador: retraso, avería, emergencia u otra situación. \\[4pt]
  Auditoría           & Registro cronológico de todas las acciones realizadas en el sistema. \\
\end{longtable}

% ════════════════════════════════════════════════════════════════════════════
\section{Preguntas Frecuentes (FAQ)}

\subsubsection*{¿Qué hago si olvidé mi contraseña?}
En la pantalla de inicio de sesión, haga clic en el enlace \textit{``¿Olvidaste tu contraseña?''}. El sistema enviará un correo con instrucciones para restablecerla.

\subsubsection*{¿Por qué no aparece una embarcación en el mapa GPS?}
Puede deberse a que su dispositivo GPS está apagado, sin señal, o fuera del rango de cobertura. Verifique el dispositivo a bordo.

\subsubsection*{¿Por qué el operador no puede iniciar el tracking?}
El Administrador debe cambiar el estado del viaje a \textbf{`En Curso'} antes de que el operador pueda seleccionar el viaje e iniciar el tracking GPS.

\subsubsection*{¿Cómo exporto un reporte en Excel?}
En el módulo de Reportes, seleccione el tipo de reporte deseado y haga clic en el botón \textbf{Excel} (verde). El archivo se descargará automáticamente.

\subsubsection*{¿Puedo eliminar un propietario que tiene embarcaciones asociadas?}
No se recomienda. Antes de eliminar un propietario, reasigne o elimine sus embarcaciones. El sistema puede impedir la eliminación si detecta registros activos relacionados.

\subsubsection*{¿Cómo sé si un ticket está confirmado?}
En el módulo de Compras y Tickets (o en la tabla de pasajeros del Operador), la columna \textbf{Estado} muestra \textit{Confirmado} en azul o \textit{Pendiente} en amarillo.

\subsubsection*{¿Quién puede acceder al módulo de Configuración?}
Todos los perfiles acceden a Configuración, pero el Administrador tiene opciones adicionales de seguridad (2FA, dispositivos confiables) que no están disponibles para el Operador.

\subsubsection*{¿Cómo reporto un incidente durante un viaje?}
Haga clic en \textbf{Reportar Incidente} en el menú lateral, complete el formulario con tipo, severidad y descripción, y haga clic en \textbf{Enviar reporte}.

% ════════════════════════════════════════════════════════════════════════════
\section{Contacto y Soporte Técnico}

Para soporte técnico, reportes de errores o consultas sobre el sistema GEONAVAL, comuníquese a través de los siguientes canales:

\begin{longtable}{LR}
  \toprule
  \textbf{Canal} & \textbf{Detalle} \\
  \midrule
  \endfirsthead
  \toprule
  \textbf{Canal} & \textbf{Detalle} \\
  \midrule
  \endhead
  \bottomrule
  \endfoot
  Teléfono           & +57 (4) 670-1234 \\[4pt]
  Correo electrónico & \href{mailto:info@geonaval.gov.co}{info@geonaval.gov.co} \\[4pt]
  Dirección          & Calle Principal \#10-20, Quibdó, Chocó, Colombia \\[4pt]
  Horario de atención & Lunes a Viernes: 8:00 AM -- 6:00 PM \quad|\quad Sábado: 8:00 AM -- 12:00 PM \\[4pt]
  Emergencias GPS    & Disponible 24/7 para incidentes críticos de monitoreo \\[4pt]
  URL del sistema    & \url{https://geonaval.onrender.com} \\
\end{longtable}

\vspace{1em}
\begin{notabox}
\textbf{Nota:} GEONAVAL es un Sistema Oficial autorizado por la Capitanía del Puerto. Versión 1.0.0 -- Mayo 2026. Todos los derechos reservados.\\
\copyright{} 2026 GEONAVAL -- Sistema de Control del Transporte Fluvial.
\end{notabox}

\vfill
\begin{center}
  \small Quibdó, Chocó -- Colombia \quad|\quad Versión 1.0.0 -- Mayo 2026
\end{center}

\end{document}
